\subsection{Uniform comparison of the scalar functionals}

When two scalar functionals are uniformly close, their supporting probabilities are close as well. To make this precise, define for two model references $\varphi,\psi$ in the preceding setting
\begin{equation}\label{eq:normalization-distance}
 \Delta(\varphi,\psi)
   =\sup_{G\in C(X)}|B_\varphi(G)-B_\psi(G)|.
\end{equation}
The comparison with $B_0$ in \cref{prop:reference-normalization-comparison} makes this number finite for every such pair. It is a pseudometric: symmetry, vanishing on the diagonal and the triangle inequality follow from the properties of the uniform norm. Distinct singularity types can have distance zero, as \cref{ex:normalization-distance-projective} shows.

\Needspace{12\baselineskip}
\begin{proposition}\label{prop:normalization-distance-dual}
For these model references,
\[
 \Delta(\varphi,\psi)
 =\sup_{\operatorname{Ent}_{\mathrm{d}V}(\nu)<\infty}
                       |H_\varphi(\nu)-H_\psi(\nu)|,
\]
where the supremum is over probability measures of finite entropy. The same value is obtained if \eqref{eq:normalization-distance} uses all bounded measurable $G$, and it does not depend on the smooth positive reference volume.
\end{proposition}
\begin{proof}
By \cref{cor:prescribed-dual-entropy}, both duals are finite exactly on the probability measures of finite entropy. A uniform bound for $|B_\varphi-B_\psi|$ passes to the defining suprema, so the supremum in the statement is at most $\Delta(\varphi,\psi)$. Conversely, denote that supremum by $D$. If $\nu\in\mathcal A_\psi(G)$, then $\nu$ has finite entropy, and the dual minimum formula gives
\[
 B_\varphi(G)
 \leq H_\varphi(\nu)-\int_XG\,\mathrm{d}\nu
 \leq B_\psi(G)+D.
\]
Interchanging the references gives the opposite inequality, so $\Delta(\varphi,\psi)\leq D$.

For bounded measurable $G$, take continuous approximations with the same uniform bound converging in $L^1(\mathrm{d}V)$; the continuity theorem \cref{thm:normalization-bounded-lone}, applied to each of the two references, extends the scalar bound to $G$. Finally, a change of volume $\mathrm{d}V'=\mathrm{e}^a\mathrm{d}V$ replaces each $B(G)$ by $B(G+a)$, and translation by $a\in C(X)$ is a bijection of the space of tests, so the supremum is unchanged.
\end{proof}

\Needspace{8\baselineskip}
\begin{proposition}\label{prop:normalization-distance-stability}\label{eq:geometric-adjoint-gap}
Let $\varphi_j,\varphi$ be model references in the same background with $\Delta(\varphi_j,\varphi)\to0$. Let $G_j,G$ be uniformly bounded measurable functions with $G_j\to G$ in $L^1(\mathrm{d}V)$, and choose $\nu_j\in\mathcal A_{\varphi_j}(G_j)$. Then the densities of the $\nu_j$ are uniformly integrable, every weak cluster point of $(\nu_j)$ belongs to $\mathcal A_\varphi(G)$, and a weakly convergent subsequence also converges weakly in $L^1$. If $\mathcal A_\varphi(G)$ is a singleton, the whole sequence converges to it in both senses.
\end{proposition}
\begin{proof}
Write $\delta_j=\Delta(\varphi_j,\varphi)$. The bounds $|B_{\varphi_j}-B_\varphi|\leq\delta_j$ and $|H_{\varphi_j}-H_\varphi|\leq\delta_j$, together with the identity $H_{\varphi_j}(\nu_j)=B_{\varphi_j}(G_j)+\int_XG_j\,\mathrm{d}\nu_j$, give
\begin{equation}\label{eq:normalization-distance-gap}
 0\leq H_\varphi(\nu_j)-\int_XG_j\,\mathrm{d}\nu_j-B_\varphi(G_j)
       \leq2\delta_j.
\end{equation}
If $\|G_j\|_\infty\leq M$, then $H_\varphi(\nu_j)\leq B_\varphi(0)+2M+2\delta_j$, so all the measures lie in one finite sublevel set of $H_\varphi$. By \cref{thm:normalization-dual-compactness}, their densities are uniformly integrable, and a weakly convergent subsequence also converges weakly in $L^1$.

Let $\nu$ be the limit of such a subsequence. Uniform integrability, the common bound and convergence in measure give $\int_X|G_j-G|\,\mathrm{d}\nu_j\to0$; weak $L^1$ convergence applies to the bounded function $G$, and \cref{thm:normalization-bounded-lone} gives $B_\varphi(G_j)\to B_\varphi(G)$. Lower semicontinuity of $H_\varphi$ in \eqref{eq:normalization-distance-gap} yields
\[
 H_\varphi(\nu)-\int_XG\,\mathrm{d}\nu-B_\varphi(G)\leq0.
\]
The minimum formula for bounded measurable densities in \cref{thm:normalization-measurable-adjoint} gives the reverse inequality, and equality means exactly that $\nu\in\mathcal A_\varphi(G)$. When this set is a singleton, compactness gives convergence of the whole sequence.
\end{proof}

The smaller-background comparison gives an explicit family of references converging to the zero reference.
